\chapter{Weighted Bipartite Matching and the Hungarian Algorithm}

The previous chapter asked a pure cardinality question about bipartite graphs: how many disjoint pairs can we form? Now we attach a number to every possible pairing and ask for the \emph{best} set of disjoint pairs. Workers and jobs, taxis and riders, antennas and frequencies -- whenever each assignment carries a cost or a profit and each side must be used at most once, you are looking at the \emph{assignment problem}. The algorithm that solves it, the Hungarian method of Kuhn and Munkres, is a standard example of \emph{primal--dual} reasoning in combinatorial optimization. It maintains a feasible solution to a dual (a set of \emph{potentials}), and uses that dual both to certify optimality and to drive the search. It is the method behind successive shortest paths for min-cost flow, specialized to bipartite graphs, where it is simple enough to run by hand.

\chapterline

\section{The Assignment Problem}
\label{ch16:assignment}

\begin{focusbox}
We are given a complete weighted bipartite graph on two sides $L$ and $R$ with $\abs{L}=\abs{R}=n$, and a cost $c(i,j)\in\R$ for assigning $i\in L$ to $j\in R$. A \emph{perfect matching} (or \emph{assignment}) is a bijection $\pi:L\to R$. We want the assignment of minimum total cost
\[
  \min_{\pi}\ \sum_{i\in L} c\bigl(i,\pi(i)\bigr).
\]
Equivalently, choosing exactly one entry from each row and each column of an $n\times n$ cost matrix $C$, minimize the sum of chosen entries.
\end{focusbox}

Two framings travel together throughout this chapter, and it pays to keep both in view.

\begin{itemize}
  \item \textbf{Graph view.} A bipartite graph $\Gph=(L\cup R, E)$ with edge weights. We seek a maximum-weight (or minimum-cost) perfect matching. This is the view in which \emph{augmenting paths} live.
  \item \textbf{Matrix view.} An $n\times n$ matrix $C$, from which we pick a \emph{transversal}: one cell per row, one per column. This is the view in which the \emph{row/column reductions} of the original Hungarian method live.
\end{itemize}

\begin{intuitionbox}
Why not just try every assignment? There are $n!$ of them. For $n=15$ that is already more than a trillion. Greedy -- repeatedly grab the cheapest available cell -- fails badly: committing $i$ to its personal favourite $j$ can strand a second worker who had \emph{only} $j$ as a cheap option, forcing an expensive repair later. The difficulty is that a locally cheap choice can be globally expensive. We need a way to make commitments that we can later \emph{revoke for free}. That is exactly what augmenting paths and dual potentials give us.
\end{intuitionbox}

\paragraph{Max-weight vs.\ min-cost; non-square inputs.} Maximizing total weight $\sum w(i,\pi(i))$ and minimizing total cost $\sum c(i,\pi(i))$ are the same problem under $c=W-w$ for any large constant $W$ (a perfect matching uses exactly $n$ edges, so shifting every weight by $W$ shifts every assignment's total by $nW$). If the two sides differ in size, pad the smaller side with dummy vertices joined by zero-cost edges; if a perfect matching need not exist or some pairings are forbidden, set forbidden costs to $+\infty$ (or a number larger than any achievable finite total). We therefore lose no generality assuming a \emph{square, complete, finite-cost} instance, and we phrase the core algorithm as minimization.

\begin{historybox}
Harold Kuhn published the algorithm in 1955 and named it the ``Hungarian method'' to honor mathematicians D\'enes K\H{o}nig and Jen\H{o} Egerv\'ary. In 1931, Egerv\'ary proved that the maximum weight of a bipartite matching equals the minimum sum of vertex weights that cover the edges -- a result that anticipates strong linear programming duality by two decades. James Munkres reorganized Kuhn's matrix version in 1957 into the standard implementation used today. In 2006 it came to light that Carl Gustav Jacobi had formulated an essentially identical method in an 1840s Latin manuscript on systems of differential equations, where finding an optimal transversal determined the order of a differential system, though his work lay largely forgotten for over a century. Kuhn's algorithm led to the general \emph{primal--dual method} in combinatorial optimization, inspiring Lester Ford and Delbert Fulkerson's network flow algorithms and Jack Edmonds's blossom algorithm. Even today, closing the complexity gap between bipartite matching and general min-cost flow for dense matrices remains an open research problem.
\end{historybox}

\section{Picture First: Tight Edges and Potentials}
\label{ch16:potentials}

Assign to each left vertex $i$ a number $u_i$ and to each right vertex $j$ a number $v_j$. Think of $u_i$ as ``how much $i$ has already been promised'' and $v_j$ as ``how much $j$ has already been promised.'' Call the pair $(u,v)$ a \emph{feasible potential} (or feasible \emph{labeling}) if it never over-promises against the cost of any edge:
\[
  u_i + v_j \ \le\ c(i,j)\qquad\text{for every } i\in L,\ j\in R.
\]
An edge $(i,j)$ is \emph{tight} (or \emph{admissible}) when its slack is zero:
\[
  \mathrm{slack}(i,j) \;=\; c(i,j) - u_i - v_j \;=\; 0 .
\]

\begin{intuitionbox}
Feasibility says: the potentials together can ``claim'' no more than the edge actually costs. The \emph{equality subgraph} $\Gph_=$ consists of exactly the tight edges -- those where the claim is exact, with no slack left over. The strategy is: \emph{find a perfect matching using only tight edges.} If you can, the matching is optimal, because the potentials hand you a matching certificate of optimality. When you cannot, you nudge the potentials -- carefully, keeping feasibility -- so that at least one new edge becomes tight, growing $\Gph_=$ until a perfect matching appears inside it.
\end{intuitionbox}

Here is the lower bound that makes potentials useful. Take \emph{any} perfect matching $\pi$ and \emph{any} feasible potential $(u,v)$. Summing the feasibility inequality over the $n$ matched edges, and using that $\pi$ hits each left vertex once and each right vertex once,
\[
  \sum_{i} c\bigl(i,\pi(i)\bigr)
  \;\ge\; \sum_i \bigl(u_i + v_{\pi(i)}\bigr)
  \;=\; \sum_{i\in L} u_i + \sum_{j\in R} v_j .
\]
The right-hand side does not depend on $\pi$. So:

\begin{whybox}
\textbf{Weak duality.} For every feasible potential $(u,v)$ and every perfect matching $\pi$,
\[
  \underbrace{\sum_{i} c\bigl(i,\pi(i)\bigr)}_{\text{primal: cost of an assignment}}
  \ \ge\
  \underbrace{\sum_i u_i + \sum_j v_j}_{\text{dual: total potential}}.
\]
Hence \emph{any} feasible potential lower-bounds the optimum, and \emph{any} matching upper-bounds it. If we ever produce a matching and a potential with \emph{equal} value, both are optimal simultaneously, and each certifies the other. Equality forces every matched edge to satisfy $u_i+v_{\pi(i)}=c(i,\pi(i))$ -- every matched edge has zero slack. This condition is \emph{complementary slackness}: whenever a primal variable $x_{ij} = 1$ is active, the corresponding dual constraint $u_i + v_j \le c(i,j)$ must be tight. The algorithm preserves this tightness invariant on matched edges while growing the matching until it is perfect.
\end{whybox}

\begin{figure}[H]
\centering
\begin{tikzpicture}[
  every node/.style={font=\small},
  Lv/.style={circle,draw=sourceblue,thick,fill=blue!10,minimum size=7.5mm,inner sep=0pt},
  Rv/.style={circle,draw=targetgreen,thick,fill=green!10,minimum size=7.5mm,inner sep=0pt},
  tight/.style={very thick,structurepurple},
  slack/.style={edgegray!60,thin,dotted},
  matched/.style={line width=1.8pt,trapred},
  lane/.style={draw=black!15,fill=black!2,rounded corners=4pt}
]
% Panel (a): Potentials and All Edges with Slacks
\begin{scope}
  \node[font=\footnotesize\bfseries,align=center] at (1.8,3.7) {\textbf{(a)} Cost Network with Potentials $(u,v)$};

  \node[Lv] (a) at (0,2.8) {$1$};
  \node[Lv] (b) at (0,1.4) {$2$};
  \node[Lv] (c) at (0,0.0) {$3$};

  \node[Rv] (x) at (3.6,2.8) {$A$};
  \node[Rv] (y) at (3.6,1.4) {$B$};
  \node[Rv] (z) at (3.6,0.0) {$C$};

  \node[left=1.5mm of a,font=\footnotesize,sourceblue] {$u_1=1$};
  \node[left=1.5mm of b,font=\footnotesize,sourceblue] {$u_2=2$};
  \node[left=1.5mm of c,font=\footnotesize,sourceblue] {$u_3=0$};

  \node[right=1.5mm of x,font=\footnotesize,targetgreen] {$v_A=3$};
  \node[right=1.5mm of y,font=\footnotesize,targetgreen] {$v_B=1$};
  \node[right=1.5mm of z,font=\footnotesize,targetgreen] {$v_C=2$};

  % Tight edges
  \draw[matched] (a) -- (x) node[pos=0.3,above,sloped,font=\scriptsize,trapred] {tight};
  \draw[tight]   (a) -- (y);
  \draw[tight]   (b) -- (y);
  \draw[matched] (c) -- (z) node[pos=0.3,below,sloped,font=\scriptsize,trapred] {tight};

  % Slack edges
  \draw[slack] (b) -- (x);
  \draw[slack] (b) -- (z);
  \draw[slack] (c) -- (x);
  \draw[slack] (c) -- (y);
  \draw[slack] (a) -- (z);

  \node[font=\footnotesize,align=center] at (1.8,-1.0)
    {Feasibility: $u_i + v_j \le c(i,j)$\\
     Slack: $c(i,j) - u_i - v_j \ge 0$};
\end{scope}

% Panel (b): Induced Equality Subgraph G_=
\begin{scope}[xshift=7.5cm]
  \node[font=\footnotesize\bfseries,align=center] at (1.8,3.7) {\textbf{(b)} Induced Equality Subgraph $\Gph_=$};

  \node[Lv] (ba) at (0,2.8) {$1$};
  \node[Lv] (bb) at (0,1.4) {$2$};
  \node[Lv] (bc) at (0,0.0) {$3$};

  \node[Rv] (bx) at (3.6,2.8) {$A$};
  \node[Rv] (by) at (3.6,1.4) {$B$};
  \node[Rv] (bz) at (3.6,0.0) {$C$};

  \node[left=1.5mm of ba,font=\footnotesize,sourceblue] {$u_1=1$};
  \node[left=1.5mm of bb,font=\footnotesize,sourceblue] {$u_2=2$};
  \node[left=1.5mm of bc,font=\footnotesize,sourceblue] {$u_3=0$};

  \node[right=1.5mm of bx,font=\footnotesize,targetgreen] {$v_A=3$};
  \node[right=1.5mm of by,font=\footnotesize,targetgreen] {$v_B=1$};
  \node[right=1.5mm of bz,font=\footnotesize,targetgreen] {$v_C=2$};

  % Only tight edges exist here
  \draw[matched] (ba) -- (bx) node[midway,above,sloped,font=\scriptsize,trapred] {matched};
  \draw[tight]   (ba) -- (by);
  \draw[tight]   (bb) -- (by);
  \draw[matched] (bc) -- (bz) node[midway,below,sloped,font=\scriptsize,trapred] {matched};

  \node[font=\footnotesize,align=center] at (1.8,-1.0)
    {Only tight edges ($\mathrm{slack}=0$) retained;\\
     free vertex $2$ needs an augmenting path};
\end{scope}
\end{tikzpicture}
\caption{The role of vertex potentials and the equality subgraph. (a)~Feasible
potentials $(u_i, v_j)$ satisfy $u_i + v_j \le c(i,j)$ across all edges. Tight
edges (purple/red) satisfy $u_i + v_j = c(i,j)$ with zero slack, while dotted
edges have strictly positive slack. (b)~The equality subgraph $\Gph_=$ discards
all positive-slack edges. Matched edges (thick red) satisfy complementary slackness.
Any augmenting path in $\Gph_=$ immediately preserves dual feasibility while
growing the cardinality of the matching.}
\label{ch16:fig-equality-subgraph}
\end{figure}

In the picture the matching $\{(1,A),(3,C)\}$ lives among tight edges, and $2$ still needs a tight option to complete a perfect matching; the algorithm's job is to reshape potentials until vertex $2$ acquires a tight edge that finishes the assignment. We now describe the algorithm that does this.

\section{Augmenting Paths in the Equality Subgraph}
\label{ch16:augment}

Recall the matching vocabulary from the unweighted chapter. Given a matching $M$, a vertex is \emph{free} if unmatched. An \emph{alternating path} alternates between non-matching and matching edges; an \emph{augmenting path} is an alternating path between two free vertices. Flipping (XOR-ing) $M$ along an augmenting path increases $\abs{M}$ by one. The new wrinkle: we only ever walk along \emph{tight} edges, so that whatever matching we build automatically satisfies complementary slackness.

\begin{focusbox}
\textbf{Invariant maintained throughout.}
\begin{enumerate}
  \item $(u,v)$ is a feasible potential: $u_i+v_j\le c(i,j)$ for all $i,j$.
  \item $M$ is a matching, and \emph{every} edge of $M$ is tight ($u_i+v_j=c(i,j)$).
\end{enumerate}
When $M$ becomes perfect, invariant (2) is exactly complementary slackness, and weak duality (Section~\ref{ch16:potentials}) makes $M$ optimal. The algorithm's only moves are: (a) augment along a tight alternating path, growing $M$; (b) adjust potentials to add a tight edge without breaking feasibility.
\end{focusbox}

\subsection{Growing a Hungarian tree}

To match one more left vertex $r$ (currently free), grow an alternating \emph{tree} rooted at $r$ inside the equality subgraph, exactly as in unweighted bipartite matching. Maintain two sets:
\begin{itemize}
  \item $S\subseteq L$: left vertices reached by the tree (including the root $r$);
  \item $T\subseteq R$: right vertices reached by the tree.
\end{itemize}
A right vertex enters $T$ via a tight edge from $S$; once in $T$, if it is matched, its matched partner enters $S$ (the only way to leave a matched right vertex along an alternating path is through its matching edge). If some right vertex in our reach is \emph{free}, we have found an augmenting path -- flip it and we are done with this phase.

The danger: every tight edge out of $S$ might land in $T$ already. Then the tree is stuck inside the current equality subgraph. This is where the dual potentials come in.

\subsection{The dual adjustment}

Suppose the tree is stuck: there is no tight edge from $S$ to $R\setminus T$. Compute the smallest slack across that frontier,
\[
  \delta \;=\; \min_{\,i\in S,\ j\in R\setminus T}\ \bigl(c(i,j) - u_i - v_j\bigr) \;>\;0 ,
\]
and update
\[
  u_i \mathrel{+}= \delta \ \ (i\in S),\qquad
  v_j \mathrel{-}= \delta \ \ (j\in T).
\]

\begin{whybox}
Why this preserves feasibility \emph{and} progress -- check each edge class $(i,j)$:
\begin{itemize}
  \item $i\in S,\ j\in T$: $u_i$ rises by $\delta$, $v_j$ falls by $\delta$; slack unchanged. \emph{Tree edges and matched edges stay tight.}
  \item $i\in S,\ j\notin T$: $u_i$ rises by $\delta$, $v_j$ fixed; slack drops by $\delta$ and remains $\ge0$ by choice of $\delta$. \emph{At least one frontier edge becomes newly tight.}
  \item $i\notin S,\ j\in T$: $u_i$ fixed, $v_j$ falls; slack rises. Still feasible.
  \item $i\notin S,\ j\notin T$: untouched.
\end{itemize}
So feasibility holds, all tree/matched edges stay tight (the tree survives intact), and the frontier gains at least one tight edge -- letting the tree grow by at least one more right vertex. After at most $n$ such adjustments the tree must reach a free right vertex, and we augment.
\end{whybox}

\begin{trapbox}
Three classic bugs.
\begin{enumerate}
  \item \textbf{Updating the wrong side of $T$.} You add $\delta$ to $u$ on $S$ and subtract from $v$ on $T$. Subtracting on all of $R$, or adding to $u$ on all of $L$, breaks feasibility for edges with one endpoint outside the tree.
  \item \textbf{Recomputing $\delta$ from scratch every time} -- correct, but $\Oh(n^2)$ per adjustment, giving $\Oh(n^4)$ overall. Maintain a \emph{slack array} $\mathrm{slk}[j]=\min_{i\in S}\mathrm{slack}(i,j)$ for $j\notin T$ and update it incrementally; this is the difference between $\Oh(n^4)$ and $\Oh(n^3)$ (Section~\ref{ch16:complexity}).
  \item \textbf{Forgetting that $\delta>0$ strictly.} If a zero-slack frontier edge existed, the tree was not actually stuck -- you should have grown it first. Mixing these up loops forever.
\end{enumerate}
\end{trapbox}

\begin{recoverybox}
\textbf{Pause and think: what happens to the dual objective during an adjustment?}
The sum of potentials changes by $\sum_{i\in S}\delta - \sum_{j\in T}\delta = (\abs{S}-\abs{T})\delta$. Because the root $r\in S$ is unmatched and every other vertex in $S$ is matched to a vertex in $T$, we have $\abs{S} = \abs{T}+1$. Thus the total dual objective strictly increases by $\delta > 0$ at every adjustment, driving the potentials steadily toward the minimum matching cost.
\end{recoverybox}

Figure~\ref{ch16:fig-dual-adjustment} shows the dual adjustment. Before the step, the Hungarian alternating tree spans $S$ on the left and $T$ on the right; the frontier edges from $S$ to $R\setminus T$ all carry positive slack, the smallest of which is $\delta$. Raising every $u_i$ for $i\in S$ by $\delta$ and lowering every $v_j$ for $j\in T$ by $\delta$ keeps every tree edge at zero slack while the frontier slacks drop by $\delta$, so at least one frontier edge enters the equality subgraph.

\begin{figure}[H]
\centering
\begin{tikzpicture}[
  every node/.style={font=\small},
  Lv/.style={circle,draw=sourceblue,thick,fill=blue!10,minimum size=7mm,inner sep=0pt},
  Rv/.style={circle,draw=targetgreen,thick,fill=green!10,minimum size=7mm,inner sep=0pt},
  Rf/.style={circle,draw=trapred,very thick,fill=red!10,minimum size=7mm,inner sep=0pt},
  tree/.style={line width=1.5pt,structurepurple},
  matched/.style={line width=1.8pt,trapred},
  front/.style={edgegray,thick,dashed},
  boxS/.style={draw=sourceblue!60,fill=blue!4,rounded corners=4pt,thick},
  boxT/.style={draw=targetgreen!60,fill=green!4,rounded corners=4pt,thick},
  boxOut/.style={draw=black!25,fill=black!2,rounded corners=4pt,dashed}
]
  % Partition boxes on Left
  \draw[boxS] (-0.7,1.1) rectangle (0.7,3.6);
  \node[sourceblue,font=\footnotesize\bfseries,align=right,anchor=east] at (-0.9,2.35) {$S \subseteq L$\\[1pt]\scriptsize($u_i \leftarrow u_i + \delta$)};

  \draw[boxOut] (-0.7,-1.2) rectangle (0.7,0.3);
  \node[edgegray,font=\footnotesize\bfseries,align=right,anchor=east] at (-0.9,-0.45) {$L \setminus S$\\[1pt]\scriptsize($u_i$ fixed)};

  % Partition boxes on Right
  \draw[boxT] (5.3,1.1) rectangle (6.7,3.6);
  \node[targetgreen,font=\footnotesize\bfseries,align=left,anchor=west] at (6.9,2.35) {$T \subseteq R$\\[1pt]\scriptsize($v_j \leftarrow v_j - \delta$)};

  \draw[boxOut] (5.3,-1.2) rectangle (6.7,0.3);
  \node[edgegray,font=\footnotesize\bfseries,align=left,anchor=west] at (6.9,-0.45) {$R \setminus T$\\[1pt]\scriptsize($v_j$ fixed)};

  % Vertices
  \node[Lv] (s1) at (0,3.0) {$s_1$};
  \node[Lv] (s2) at (0,1.7) {$s_2$};
  \node[Lv,fill=black!5,draw=edgegray] (l_out) at (0,-0.45) {$u_{\mathrm{out}}$};

  \node[Rv] (t1) at (6.0,3.0) {$t_1$};
  \node[Rv] (t2) at (6.0,1.7) {$t_2$};
  \node[Rf] (f1) at (6.0,-0.45) {$r_{\mathrm{free}}$};

  % Edges between S and T (Tree and Matched)
  \draw[tree]    (s1) -- (t1) node[midway,above=2pt,font=\scriptsize,structurepurple] {tree edge ($\Delta \mathrm{slack} = 0$)};
  \draw[matched] (s2) -- (t2) node[midway,above=2pt,font=\scriptsize,trapred] {matched edge ($\Delta \mathrm{slack} = 0$)};

  % Frontier edge from S to R \ T
  \draw[front] (s2) -- (f1) node[pos=0.74,sloped,above=2pt,font=\scriptsize,edgegray] {slack $\delta \to 0$};

  % Edge from L \ S to T
  \draw[dotted,thick,edgegray] (l_out) -- (t2) node[pos=0.26,sloped,below=2pt,font=\scriptsize,edgegray] {slack $+\delta$};

  % Legend / Rule box below
  \node[draw=black!20,rounded corners=3pt,fill=white,font=\scriptsize,align=center,inner sep=4pt] at (3.0,-2.35) {
    \textbf{Slack Updates by Edge Class:}\\
    $S \to T$: $\Delta = (+\delta) - (+\delta) = 0$ (all tree and matched edges remain tight)\\
    $S \to R \setminus T$: $\Delta = (+\delta) - (0) \implies \mathrm{slack}' = \mathrm{slack} - \delta \ge 0$ (frontier tightens; $\delta = \min$ hits 0)\\
    $L \setminus S \to T$: $\Delta = (0) - (-\delta) \implies \mathrm{slack}' = \mathrm{slack} + \delta \ge 0$ (feasibility preserved)\\
    $L \setminus S \to R \setminus T$: $\Delta = 0$ (completely unaffected)
  };
\end{tikzpicture}
\caption{Dual adjustment when the Hungarian tree is stuck.
When no tight edge crosses from $S$ to $R\setminus T$, potentials are shifted by
the minimum frontier slack $\delta = \min_{i\in S, j\notin T}(c(i,j)-u_i-v_j) > 0$.
Increasing $u_i$ by $\delta$ on $S$ and decreasing $v_j$ by $\delta$ on $T$ preserves
exact zero slack for all internal tree and matched edges ($S \to T$), while decreasing
frontier slacks ($S \to R\setminus T$) by $\delta$. The edge achieving $\delta$ becomes
tight, expanding the alternating tree without violating feasibility anywhere.}
\label{ch16:fig-dual-adjustment}
\end{figure}

\begin{walkbox}
Why must $\delta$ be finite -- why is there always a frontier edge to measure? Because the instance is complete: every $i\in S$ has an edge to every $j\in R\setminus T$, and as long as the matching is not yet perfect there is at least one right vertex outside $T$ (otherwise $T=R$ would already contain a free right vertex, and we would have augmented instead of adjusting). So the minimum in the definition of $\delta$ ranges over a nonempty set and is a genuine finite number. On a \emph{sparse} instance you must instead treat missing edges as cost $+\infty$, and then guard against an all-infinite frontier, which signals that no perfect matching exists.
\end{walkbox}

\section{The Hungarian Algorithm}
\label{ch16:algorithm}

We give the $\Oh(n^3)$ form directly: match left vertices one at a time, and within each phase maintain the slack array so each phase costs $\Oh(n^2)$. The implementation below uses the standard ``virtual root'' / column-driven Jonker--Volgenant style augmentation, which is both fast and short. We write $\mathrm{matchR}[j]$ for the left vertex matched to right vertex $j$ ($0$ if free), and reconstruct the path with a parent array $\mathrm{prev}[j]$.

\begin{algorithm}[H]
\DontPrintSemicolon
\caption{Hungarian algorithm (min-cost perfect assignment, $\Oh(n^3)$)}
\KwIn{$n\times n$ cost matrix $c$ (1-indexed); rows $L=\set{1..n}$, columns $R=\set{1..n}$.}
\KwOut{potentials $u,v$ and assignment $\mathrm{matchR}$ of minimum total cost.}
$u[0..n]\gets 0$;\quad $v[0..n]\gets 0$;\quad $\mathrm{matchR}[0..n]\gets 0$\;
\For{$r\gets 1$ \KwTo $n$}{
  $j_0\gets 0$;\quad $\mathrm{matchR}[0]\gets r$\tcp*{column $0$ is the virtual root carrying new row $r$}
  $\mathrm{slk}[j]\gets +\infty,\ \mathrm{prev}[j]\gets 0$ for all $j\in R$;\quad $\mathrm{used}[0..n]\gets \textsc{false}$\;
  \Repeat{$\mathrm{matchR}[j_0]=0$}{
    $\mathrm{used}[j_0]\gets\textsc{true}$;\quad $i\gets \mathrm{matchR}[j_0]$;\quad $\delta\gets+\infty$;\quad $j_1\gets -1$\;
    \For{$j\in R$ with $\mathrm{used}[j]=\textsc{false}$}{
      $\mathrm{cur}\gets c[i][j]-u[i]-v[j]$\tcp*{current slack of edge $(i,j)$}
      \If{$\mathrm{cur}<\mathrm{slk}[j]$}{$\mathrm{slk}[j]\gets \mathrm{cur}$;\quad $\mathrm{prev}[j]\gets j_0$}
      \If{$\mathrm{slk}[j]<\delta$}{$\delta\gets \mathrm{slk}[j]$;\quad $j_1\gets j$}
    }
    \For{$j\in \set{0..n}$}{
      \eIf{$\mathrm{used}[j]$}{$u[\mathrm{matchR}[j]]\mathrel{+}=\delta$;\quad $v[j]\mathrel{-}=\delta$}
                          {$\mathrm{slk}[j]\mathrel{-}=\delta$}
    }
    $j_0\gets j_1$\tcp*{step to the newly tightened column}
  }
  \While{$j_0\ne 0$}{
    $j_1\gets \mathrm{prev}[j_0]$;\quad $\mathrm{matchR}[j_0]\gets \mathrm{matchR}[j_1]$;\quad $j_0\gets j_1$\tcp*{flip the path}
  }
}
\Return{$u,v,\mathrm{matchR}$}\;
\end{algorithm}

\begin{intuitionbox}
Read the inner \textbf{Repeat} as ``grow the Hungarian tree by one column per iteration.'' The virtual column $0$ holds the freshly added row $r$. Each iteration: scan slacks of the current row $i=\mathrm{matchR}[j_0]$ into the slack array, find the minimum residual slack $\delta$ over unused columns, push \emph{all} potentials by $\delta$ (this is the dual adjustment, applied lazily through the slack array), and step to the newly tightened column $j_1$. The loop stops the moment it steps onto a \emph{free} column ($\mathrm{matchR}[j_1]=0$); then the \textbf{While} walks $\mathrm{prev}$ back to the root, flipping matched/unmatched edges -- a complete augmentation. Notice potentials are updated \emph{inside} the search, so they stay consistent with the partial tree.
\end{intuitionbox}

\begin{whybox}
Why the lazy potential update is correct. We maintain $\mathrm{slk}[j]=\min_{i\in S}\mathrm{slack}(i,j)$ for $j\notin T$, where $S$ is the set of rows pulled into the tree and $T$ the set of used columns. A dual adjustment of size $\delta$ must (i) leave tree columns ($\mathrm{used}[j]$, i.e.\ $j\in T$) tight, and (ii) reduce every outside column's recorded slack by $\delta$. The second inner loop does exactly that: for used columns it shifts $u$ of the matched row up and $v$ down (so the slack of tree edges is unchanged), and for unused columns it decrements $\mathrm{slk}[j]$ by $\delta$. Whenever some $\mathrm{slk}[j]$ reaches $0$, the edge into $j$ has become tight -- and that is the $j_1$ we step to next. This is the Section~\ref{ch16:augment} adjustment executed in $\Oh(n)$ per step rather than $\Oh(n^2)$.
\end{whybox}

\section{Worked Example, Run by Hand}
\label{ch16:worked}

Take the $3\times3$ cost matrix (minimize). Rows are workers $1,2,3$; columns are jobs $A,B,C$.
\[
  C=\begin{array}{c|ccc}
      & A & B & C\\\hline
    1 & 4 & 1 & 3\\
    2 & 2 & 0 & 5\\
    3 & 3 & 2 & 2
  \end{array}
\]

\begin{labbox}
Before running the full algorithm, get the answer by brute force so we have something to check against. The six assignments and their costs:
\[
\begin{aligned}
&1A,2B,3C:\ 4+0+2=6 &\quad& 1A,2C,3B:\ 4+5+2=11\\
&1B,2A,3C:\ 1+2+2=5 &\quad& 1B,2C,3A:\ 1+5+3=9\\
&1C,2A,3B:\ 3+2+2=7 &\quad& 1C,2B,3A:\ 3+0+3=6
\end{aligned}
\]
The minimum is $5$, achieved by $1\!\to\!B,\ 2\!\to\!A,\ 3\!\to\!C$. Keep $5$ in your pocket.
\end{labbox}

Now run the Hungarian method in the dual-potential form. It is easiest to track by hand using the classic ``reduce then cover'' presentation, which is the same algorithm: subtracting the row minimum from a row decreases $c$ by a constant on that row, i.e.\ raises $u_i$; subtracting a column minimum raises $v_j$. Tight edges are the zeros.

\textbf{Step 1 -- row reduction (initialize $u_i$ to the row minimum).}
Row minima are $u_1=1,\ u_2=0,\ u_3=2$. Subtract:
\[
C^{(1)}=\begin{array}{c|ccc}
   & A & B & C\\\hline
 1 & 3 & \mathbf{0} & 2\\
 2 & 2 & \mathbf{0} & 5\\
 3 & 1 & \mathbf{0} & \mathbf{0}
\end{array}
\qquad u=(1,0,2),\ v=(0,0,0).
\]
Bold zeros are tight edges. Note all three rows put their zero in $B$ -- a collision.

\textbf{Step 2 -- column reduction (initialize $v_j$).}
Column minima of $C^{(1)}$: $v_A=1,\ v_B=0,\ v_C=0$. Subtract from column $A$:
\[
C^{(2)}=\begin{array}{c|ccc}
   & A & B & C\\\hline
 1 & 2 & \mathbf{0} & 2\\
 2 & 1 & \mathbf{0} & 5\\
 3 & \mathbf{0} & \mathbf{0} & \mathbf{0}
\end{array}
\qquad u=(1,0,2),\ v=(1,0,0).
\]

\textbf{Step 3 -- try to match on the zeros (equality subgraph).}
Tight edges now: $1\!-\!B$, $2\!-\!B$, $3\!-\!A$, $3\!-\!B$, $3\!-\!C$. Greedily match: take $3\!-\!A$, then $1\!-\!B$. Vertex $2$ has only $2\!-\!B$ tight, but $B$ is taken. Grow a Hungarian tree from free worker $2$:
\[
S=\{2\},\quad 2\!-\!B \text{ tight} \Rightarrow B\in T,\quad B \text{ matched to } 1 \Rightarrow 1\in S.
\]
Now $S=\{1,2\}$, $T=\{B\}$. Tight edges from $S$ to $R\setminus T=\{A,C\}$? From $1$: $1\!-\!A$ slack $2$, $1\!-\!C$ slack $2$. From $2$: $2\!-\!A$ slack $1$, $2\!-\!C$ slack $5$. None is tight -- the tree is stuck, as shown in Figure~\ref{ch16:fig-worked-stalled}, where every frontier edge out of $S=\{1,2\}$ carries positive slack and $(2,A)$ achieves the minimum $\delta=1$.

\begin{figure}[H]
\centering
\begin{tikzpicture}[
  every node/.style={font=\small},
  Lv/.style={circle,draw=sourceblue,thick,fill=blue!10,minimum size=7.5mm,inner sep=0pt},
  Rv/.style={circle,draw=targetgreen,thick,fill=green!10,minimum size=7.5mm,inner sep=0pt},
  tight/.style={very thick,structurepurple},
  matched/.style={line width=2pt,trapred},
  front/.style={edgegray,thick,dashed},
  minfront/.style={quotientorange,line width=1.8pt,dashed}
]
  \node[Lv] (w1) at (0,3.4) {$1$};
  \node[Lv] (w2) at (0,1.7) {$2$};
  \node[Lv,fill=black!5,draw=edgegray] (w3) at (0,0.0) {$3$};

  \node[Rv,fill=black!5,draw=edgegray] (jA) at (4.6,3.4) {$A$};
  \node[Rv] (jB) at (4.6,1.7) {$B$};
  \node[Rv,fill=black!5,draw=edgegray] (jC) at (4.6,0.0) {$C$};

  \node[left=2mm of w1,font=\footnotesize,sourceblue] {$u_1=1\ (S)$};
  \node[left=2mm of w2,font=\footnotesize,sourceblue] {$u_2=0\ (S)$};
  \node[left=2mm of w3,font=\footnotesize,edgegray] {$u_3=2$};

  \node[right=2mm of jA,font=\footnotesize,edgegray] {$v_A=1$};
  \node[right=2mm of jB,font=\footnotesize,targetgreen] {$v_B=0\ (T)$};
  \node[right=2mm of jC,font=\footnotesize,edgegray] {$v_C=0$};

  % Tree edges and current matching in S x T
  \draw[matched] (w1)--(jB) node[pos=0.75,above=2pt,sloped,font=\scriptsize,trapred]{matched};
  \draw[tight]   (w2)--(jB) node[midway,below=2pt,font=\scriptsize,structurepurple]{tree ($0$)};

  % Frontier edges from S to {A, C}
  \draw[front]    (w1)--(jA) node[midway,above=2pt,font=\scriptsize,edgegray]{slack $2$};
  \draw[minfront] (w2)--(jA) node[pos=0.25,above=2pt,sloped,font=\scriptsize,quotientorange]{$\delta = 1$};
  \draw[front]    (w2)--(jC) node[midway,below=2pt,sloped,font=\scriptsize,edgegray]{slack $5$};
\end{tikzpicture}
\caption{Step 3 of the worked example: the Hungarian tree rooted at free worker $2$ reaches $B$ via tight edge $(2,B)$ and pulls matched worker $1$ into $S$, stalling at $S=\{1,2\}$ and $T=\{B\}$. Among the frontier edges from $S$ to $R\setminus T=\{A,C\}$, edge $(2,A)$ has the smallest positive slack $\delta = 1$.}
\label{ch16:fig-worked-stalled}
\end{figure}

\textbf{Step 4 -- dual adjustment.}
\[
\delta=\min\{2,2,1,5\}=1\ \ (\text{edge } 2\!-\!A).
\]
Raise $u$ on $S=\{1,2\}$ by $1$, lower $v$ on $T=\{B\}$ by $1$:
\[
u=(2,1,2),\quad v=(1,-1,0).
\]
Recompute slacks (equivalently: subtract $\delta$ from uncovered cells, add it to doubly-covered cells):
\[
C^{(3)}=\begin{array}{c|ccc}
   & A & B & C\\\hline
 1 & 1 & \mathbf{0} & 1\\
 2 & \mathbf{0} & \mathbf{0} & 4\\
 3 & \mathbf{0} & 1 & \mathbf{0}
\end{array}
\]
Edge $2\!-\!A$ is now tight. The tree extends: $A\in T$ via $2\!-\!A$; $A$ is matched to $3$, so $3\in S$. From $3$, edge $3\!-\!C$ is tight and $C$ is \emph{free}. Augmenting path found:
\[
2 \;\xrightarrow{\text{free}}\; A \;\xrightarrow{\text{match}}\; 3 \;\xrightarrow{\text{free}}\; C.
\]

Figure~\ref{ch16:fig-worked-augment} contrasts the equality subgraph $\Gph_=$ right after this dual step -- showing the newly tightened edge $(2,A)$ completing the augmenting path $2\to A\to 3\to C$ -- with the optimal perfect matching obtained after flipping that path in Step~5.

\begin{figure}[H]
\centering
\begin{tikzpicture}[
  every node/.style={font=\small},
  Lv/.style={circle,draw=sourceblue,thick,fill=blue!10,minimum size=7.5mm,inner sep=0pt},
  Rv/.style={circle,draw=targetgreen,thick,fill=green!10,minimum size=7.5mm,inner sep=0pt},
  tight/.style={very thick,structurepurple},
  matched/.style={line width=2pt,trapred},
  aug/.style={line width=2pt,quotientorange,dashed},
]
% Panel (a): Before flip
\begin{scope}
  \node[font=\footnotesize\bfseries] at (1.8,3.8) {\textbf{(a)} Step 4: Augmenting Path in $\Gph_=$};

  \node[Lv] (w1) at (0,3) {$1$};
  \node[Lv] (w2) at (0,1.5) {$2$};
  \node[Lv] (w3) at (0,0) {$3$};
  \node[Rv] (jA) at (3.6,3) {$A$};
  \node[Rv] (jB) at (3.6,1.5) {$B$};
  \node[Rv] (jC) at (3.6,0) {$C$};

  \node[left=1.5mm of w1,font=\footnotesize,sourceblue] {$u_1=2$};
  \node[left=1.5mm of w2,font=\footnotesize,sourceblue] {$u_2=1$};
  \node[left=1.5mm of w3,font=\footnotesize,sourceblue] {$u_3=2$};

  \node[right=1.5mm of jA,font=\footnotesize,targetgreen] {$v_A=1$};
  \node[right=1.5mm of jB,font=\footnotesize,targetgreen] {$v_B=-1$};
  \node[right=1.5mm of jC,font=\footnotesize,targetgreen] {$v_C=0$};

  % all tight edges of C^(3): 1-B, 2-A, 2-B, 3-A, 3-C
  \draw[tight]   (w1)--(jB);
  \draw[tight]   (w2)--(jB);
  \draw[tight]   (w3)--(jA);
  % matched edge before the flip: 3-A and 1-B
  \draw[matched] (w3)--(jA) node[pos=0.25,above=2pt,sloped,font=\scriptsize,trapred]{matched};
  \draw[matched] (w1)--(jB) node[pos=0.25,above=2pt,sloped,font=\scriptsize,trapred]{matched};
  % augmenting path edges 2-A (free), 3-C (free); 3-A is the match in the path
  \draw[aug]     (w2)--(jA) node[pos=0.75,above=2pt,sloped,font=\scriptsize,quotientorange]{aug free};
  \draw[aug]     (w3)--(jC) node[pos=0.5,below=2pt,sloped,font=\scriptsize,quotientorange]{aug free};

  \node[font=\footnotesize,align=center] at (1.8,-0.9)
    {Augmenting path $2 \to A \to 3 \to C$};
\end{scope}

% Panel (b): After flip
\begin{scope}[xshift=7.2cm]
  \node[font=\footnotesize\bfseries] at (1.8,3.8) {\textbf{(b)} Step 5: Optimal Assignment After Flip};

  \node[Lv] (bw1) at (0,3) {$1$};
  \node[Lv] (bw2) at (0,1.5) {$2$};
  \node[Lv] (bw3) at (0,0) {$3$};
  \node[Rv] (bjA) at (3.6,3) {$A$};
  \node[Rv] (bjB) at (3.6,1.5) {$B$};
  \node[Rv] (bjC) at (3.6,0) {$C$};

  \node[left=1.5mm of bw1,font=\footnotesize,sourceblue] {$u_1=2$};
  \node[left=1.5mm of bw2,font=\footnotesize,sourceblue] {$u_2=1$};
  \node[left=1.5mm of bw3,font=\footnotesize,sourceblue] {$u_3=2$};

  \node[right=1.5mm of bjA,font=\footnotesize,targetgreen] {$v_A=1$};
  \node[right=1.5mm of bjB,font=\footnotesize,targetgreen] {$v_B=-1$};
  \node[right=1.5mm of bjC,font=\footnotesize,targetgreen] {$v_C=0$};

  % non-matched tight edges
  \draw[tight,opacity=0.35] (bw2)--(bjB);
  \draw[tight,opacity=0.35] (bw3)--(bjA);
  % final matched edges
  \draw[matched] (bw1)--(bjB) node[pos=0.32,above=2pt,sloped,font=\scriptsize,trapred]{$c=1$};
  \draw[matched] (bw2)--(bjA) node[pos=0.68,above=2pt,sloped,font=\scriptsize,trapred]{$c=2$};
  \draw[matched] (bw3)--(bjC) node[midway,below=2pt,font=\scriptsize,trapred]{$c=2$};

  \node[font=\footnotesize,align=center] at (1.8,-0.9)
    {Cost $1+2+2=5 = \sum u_i + \sum v_j$};
\end{scope}
\end{tikzpicture}
\caption{Completing the worked example. (a)~The equality subgraph after the $\delta=1$ dual adjustment, with augmenting path $2\to A\to 3\to C$. (b)~Flipping the path replaces $(3,A)$ with $(2,A)$ and $(3,C)$, yielding the optimal assignment of cost $5$.}
\label{ch16:fig-worked-augment}
\end{figure}

\textbf{Step 5 -- flip the path.} Currently $3\!-\!A$ is matched; flipping along $2\!-\!A\!-\!3\!-\!C$ makes $2\!-\!A$ and $3\!-\!C$ matched and frees $A$ from $3$. Keeping $1\!-\!B$, the final matching is
\[
\boxed{1\to B,\quad 2\to A,\quad 3\to C}\qquad\text{cost } 1+2+2 = 5 .
\]

\textbf{Step 6 -- read the certificate.} Final potentials $u=(2,1,2)$, $v=(1,-1,0)$. Dual value
\[
\sum_i u_i + \sum_j v_j = (2+1+2)+(1-1+0) = 5+0 = 5,
\]
matching the primal cost $5$. Weak duality with equality $\Rightarrow$ optimal. Every matched edge is tight: $c(1,B)-u_1-v_B = 1-2-(-1)=0$, $c(2,A)-u_2-v_A=2-1-1=0$, $c(3,C)-u_3-v_C=2-2-0=0$. Complementary slackness holds on the nose.

\begin{labbox}
\textbf{Trace it yourself.} Re-run on
\(
\left[\begin{smallmatrix} 7&5&11\\ 9&8&5\\ 4&6&5\end{smallmatrix}\right]
\)
(minimize). Predict: row reduction gives $u=(5,5,4)$, and here the reduced zeros already land in distinct columns, so no dual step is needed at all. Brute-force the six assignments first, then confirm your final $u,v$ sum to the optimum. (Spoiler: the optimum is $14$, via $1\!\to\!B,\,2\!\to\!C,\,3\!\to\!A$.)
\end{labbox}

\subsection{A second trace, narrated in full}

Before leaving the hand trace, let us run the matrix from the box above to see what happens when row reduction alone already isolates a full transversal of zeros -- and why small cost perturbations leave that transversal intact.
\[
  C=\begin{array}{c|ccc}
      & A & B & C\\\hline
    1 & 7 & 5 & 11\\
    2 & 9 & 8 & 5\\
    3 & 4 & 6 & 5
  \end{array}
\]
Row reduction subtracts $u=(5,5,4)$:
\[
C^{(1)}=\begin{array}{c|ccc}
   & A & B & C\\\hline
 1 & 2 & \mathbf{0} & 6\\
 2 & 4 & 3 & \mathbf{0}\\
 3 & \mathbf{0} & 2 & 1
\end{array}
\qquad u=(5,5,4),\ v=(0,0,0).
\]
Each column already contains a zero, so column reduction adds nothing: $v=(0,0,0)$. The tight edges are $1\!-\!B$, $2\!-\!C$, $3\!-\!A$ -- one per row and per column, a perfect matching on the nose. There is no collision, so the algorithm halts immediately with
\[
\boxed{1\to B,\ 2\to C,\ 3\to A},\qquad \text{cost } 5+5+4 = 14,
\]
and the dual value $\sum u_i + \sum v_j = 14 + 0 = 14$ confirms it. ($v$ is all zero here, so the row potentials alone carry the bound.) Here row reduction alone already exposed a perfect matching. Even if we raise $c(2,C)$ from $5$ to $7$, row $2$'s minimum simply shifts from $5$ to $7$:
\[
C'=\left[\begin{smallmatrix} 7&5&11\\ 9&8&7\\ 4&6&5\end{smallmatrix}\right]
\ \xrightarrow{\,u=(5,7,4)\,}\
C'^{(1)}=\begin{array}{c|ccc}
   & A & B & C\\\hline
 1 & 2 & \mathbf{0} & 6\\
 2 & 2 & 1 & \mathbf{0}\\
 3 & \mathbf{0} & 2 & 1
\end{array}
\]
Now the zeros still land in three distinct columns, so again no dual step is needed -- the assignment $1\!\to\!B,2\!\to\!C,3\!\to\!A$ stays optimal at cost $5+7+4=16$. The general point: \emph{whenever the reduced matrix has a transversal of zeros, you are already done}, and the dual adjustments exist to create such a transversal when row/column reduction alone does not. The first worked example needed one adjustment; some instances need up to $n-1$ per phase.

\begin{trapbox}
\textbf{Hand-execution pitfalls.}
\begin{enumerate}
  \item \textbf{Declaring victory at $n$ zeros.} Having $n$ (or more) zeros is necessary but not sufficient -- they must form a \emph{transversal} (no two in the same row or column). The first example had four zeros after step~2 but they shared column $B$, so no transversal existed yet.
  \item \textbf{Forgetting to update $v$ negative.} The duals on the equality constraints are free, so $v_j<0$ is normal and not a sign of error. Clamping $v$ at $0$ silently breaks the certificate.
  \item \textbf{Re-reducing rows mid-run.} After the initial reduction, you never subtract a fresh row minimum again; all further changes come through the structured $S/T$ dual step. A stray re-reduction can break invariant~(2).
\end{enumerate}
\end{trapbox}

\begin{recoverybox}
\textbf{Pause and think: why does matrix reduction never change the optimal assignment?}
Every candidate assignment chooses exactly one entry from each row $i$ and each column $j$. If you subtract $u_i$ from row $i$ and $v_j$ from column $j$, the cost of \emph{every} candidate assignment changes by the exact same constant $\sum_i u_i + \sum_j v_j$. So potential shifts do not change the ranking of assignments.
\end{recoverybox}

\section{Correctness}
\label{ch16:correctness}

\begin{theorem}[Optimality of the Hungarian algorithm]
\label{ch16:thm-correct}
On a square instance with finite costs, the algorithm terminates with a perfect matching $\pi$ and a feasible potential $(u,v)$ such that $\sum_i c(i,\pi(i)) = \sum_i u_i + \sum_j v_j$. Consequently $\pi$ is a minimum-cost perfect matching.
\end{theorem}

\begin{proof}
We verify the two invariants of Section~\ref{ch16:augment} are preserved, that each phase terminates, and that the final equality yields optimality.

\emph{Feasibility (invariant 1).} Initialization sets $u_i=\min_j c(i,j)$ and then $v_j=\min_i (c(i,j)-u_i)\ge 0$ shifts, giving $u_i+v_j\le c(i,j)$ for all $i,j$. Each dual adjustment of size $\delta=\min_{i\in S,\,j\notin T}\mathrm{slack}(i,j)$ keeps feasibility by the case analysis in the ``Why it works'' box of Section~\ref{ch16:augment}: only $S$-rows rise and only $T$-columns fall, and the binding constraints are the frontier $S\times(R\setminus T)$, whose slacks drop by exactly $\delta\ge0$ and stay nonnegative.

\emph{Tight matching (invariant 2).} Initially $M=\varnothing$ satisfies it vacuously. Dual adjustments leave all matched edges tight: a matched edge has both endpoints in the tree ($i\in S,\ j\in T$, slack unchanged). Augmentation flips along an alternating path of tight edges, so every new matched edge is tight.

\emph{Each phase terminates.} A phase grows a tree rooted at the new free vertex $r$. Each growth step adds a right vertex to $T$ (and, if it is matched, its mate to $S$); each dual adjustment makes at least one frontier edge tight, enabling a growth step. Since $\abs{T}\le n$, after at most $n$ adjustments and $n$ growth steps the tree reaches a free right vertex. That such a vertex is reachable is guaranteed by a K\H{o}nig-style counting argument: as long as no free right vertex is reachable, $T$ is exactly the set of mates of $S\setminus\{r\}$, so $\abs{T}=\abs{S}-1$; the dual step strictly enlarges the equality subgraph on the frontier, and since the instance is complete, $S$ cannot grow past $n$ without exposing a free right vertex. Augmentation then increases $\abs{M}$ by $1$.

\emph{Overall termination and optimality.} There are exactly $n$ phases, each increasing $\abs{M}$ by one, so $M$ becomes perfect. At that point invariant (2) gives $u_i+v_{\pi(i)}=c(i,\pi(i))$ for every matched edge. Summing over the $n$ matched edges,
\[
  \sum_i c(i,\pi(i)) = \sum_i (u_i+v_{\pi(i)}) = \sum_i u_i + \sum_j v_j .
\]
Combined with weak duality (Section~\ref{ch16:potentials}) -- the left side is $\ge$ the right for every matching, and the right side is $\le$ the optimum for every feasible potential -- both are optimal: no matching is cheaper and no potential is larger.
\end{proof}

\begin{whybox}
\textbf{The exchange view of optimality}, for readers who prefer it to duality. Suppose $\pi$ is the matching the algorithm returns and some other perfect matching $\pi'$ is strictly cheaper. The symmetric difference $\pi\oplus\pi'$ is a disjoint union of alternating cycles. Walk around any such cycle: swapping $\pi$-edges for $\pi'$-edges changes the cost by $\sum(\pi'\text{-edges}) - \sum(\pi\text{-edges})$. Now telescope the potentials around the cycle -- each vertex is entered and left once, so the $u,v$ contributions cancel, giving
\[
  \sum(\pi'\text{-edges}) - \sum(\pi\text{-edges}) = \sum_{e\in\pi'\cap\text{cycle}} \mathrm{slack}(e) - \sum_{e\in\pi\cap\text{cycle}} \mathrm{slack}(e).
\]
Every $\pi$-edge is tight ($\mathrm{slack}=0$) and every $\pi'$-edge has $\mathrm{slack}\ge0$, so the right side is $\ge0$: no cycle can decrease the cost. Hence $\pi$ is optimal. This is min-cost flow's ``no negative residual cycle'' criterion, specialized to matchings.
\end{whybox}

\section{Complexity}
\label{ch16:complexity}

\begin{theorem}[Running time]
\label{ch16:thm-time}
The algorithm of Section~\ref{ch16:algorithm} runs in $\Oh(n^3)$ time and $\Oh(n^2)$ space (or $\Oh(n)$ extra space beyond the input matrix).
\end{theorem}

\begin{proof}
There are $n$ phases (one per left vertex). Within a phase, the inner \textbf{Repeat} executes at most $n$ times: each iteration marks a previously-unused column $j_0$ as used and never unmarks it, so with $n$ columns the loop steps at most $n$ times before hitting a free column. Each iteration does $\Oh(n)$ work -- the slack scan over unused columns is $\Oh(n)$, and the potential-update loop over all columns is $\Oh(n)$. So a phase costs $\Oh(n^2)$, and the whole run $\Oh(n\cdot n^2)=\Oh(n^3)$. The path-flip \textbf{While} is $\Oh(n)$ per phase, dominated. Space: $u,v,\mathrm{matchR},\mathrm{slk},\mathrm{prev},\mathrm{used}$ are $\Oh(n)$ each, plus the $\Oh(n^2)$ cost matrix.
\end{proof}

\begin{complexitybox}
\begin{tabular}{ll}
\toprule
Variant & Time \\
\midrule
Naive ``reduce \& cover'' (recompute $\delta$ each step)      & $\Oh(n^4)$ \\
Hungarian with slack array (this chapter)                     & $\Oh(n^3)$ \\
Min-cost flow, $n$ shortest-path phases (dense)               & $\Oh(n(m+n\log n)) = \Oh(n^3)$ \\
Jonker--Volgenant (shortest-augmenting refinements)           & $\Oh(n^3)$ worst, faster in practice \\
Scaling (Gabow--Tarjan), integer costs in $[0,C]$             & $\Oh(\sqrt{n}\,m\,\log(nC))$ \\
\bottomrule
\end{tabular}
\\[0.4em]
Bottom line: $\Oh(n^3)$ is the standard, and is excellent up to a few thousand on each side. For sparse or huge instances, route through min-cost flow (next box) or a scaling algorithm.
\end{complexitybox}

\begin{connectionbox}
\textbf{The Hungarian algorithm as a specialization of min-cost flow} (Chapter~14). Build a flow network: source $s\to$ each $i\in L$ with capacity $1$, cost $0$; each edge $i\to j$ with capacity $1$, cost $c(i,j)$; each $j\in R\to t$ with capacity $1$, cost $0$. A min-cost integral flow of value $n$ from $s$ to $t$ is a min-cost perfect matching. The potentials $u,v$ are precisely the \emph{node potentials} of the successive-shortest-paths algorithm, and the dual adjustment is the reduced-cost reweighting (Johnson's potentials) that keeps every reduced cost nonnegative so Dijkstra applies. Each phase $=$ one shortest augmenting path. The Hungarian algorithm is the bipartite-assignment specialization where the layered structure turns the shortest-path tree into the Hungarian tree, and a Bellman--Ford-free Dijkstra runs in $\Oh(n^2)$ per phase.
\end{connectionbox}

\section{Primal--Dual View and Complementary Slackness}
\label{ch16:lp}

The assignment problem has a natural linear programming formulation that explains where the vertex potentials come from. Let $x_{ij}\in\{0,1\}$ indicate whether worker $i$ is assigned to job $j$. The integer programming formulation is:
\[
\min \sum_{i,j} c_{ij}x_{ij}
\quad\text{s.t.}\quad
\sum_j x_{ij}=1\ (\forall i\in L),\quad \sum_i x_{ij}=1\ (\forall j\in R),\quad x_{ij}\ge 0 .
\]
The constraint matrix is the vertex-edge incidence matrix of the complete bipartite graph $K_{n,n}$. Because $K_{n,n}$ is bipartite, this matrix is \emph{totally unimodular} (every square submatrix has determinant in $\set{-1, 0, 1}$). Consequently, the polyhedron defined by the linear relaxation $x_{ij} \ge 0$ has integer vertices, ensuring that an optimal basic feasible solution is automatically an integer matching vector $x \in \{0,1\}^{n\times n}$. Continuous LP algorithms and combinatorial matching algorithms optimize over the exact same convex hull.

\begin{focusbox}
\textbf{Primal LP (P)} \hfill \textbf{Dual LP (D)}
\[
\begin{aligned}
\min\ & \textstyle\sum_{ij} c_{ij}x_{ij}\\
& \textstyle\sum_j x_{ij}=1 && (u_i)\\
& \textstyle\sum_i x_{ij}=1 && (v_j)\\
& x_{ij}\ge 0
\end{aligned}
\qquad\qquad
\begin{aligned}
\max\ & \textstyle\sum_i u_i + \sum_j v_j\\
& u_i + v_j \le c_{ij} && (x_{ij})\\
& u_i, v_j\ \text{free}
\end{aligned}
\]
\end{focusbox}

The dual variables are exactly our potentials, and the dual constraint $u_i+v_j\le c_{ij}$ is exactly feasibility. The equality-constraint duals are free (unsigned), which is why $v_B$ went negative in the worked example without alarm. LP duality gives \emph{strong} duality: $\max(\mathrm{D})=\min(\mathrm{P})$. The Hungarian algorithm is a constructive proof of strong duality for this specific LP.

\begin{whybox}
\textbf{Complementary slackness}, stated for this pair. A primal-feasible $x$ and dual-feasible $(u,v)$ are simultaneously optimal iff for every edge
\[
  x_{ij} > 0 \ \Longrightarrow\ u_i+v_j = c_{ij}
  \qquad(\text{``every used edge is tight''}).
\]
(The dual-slack $\times$ primal-slack condition on the constraints is automatic, since the primal constraints are equalities.) The algorithm maintains dual feasibility \emph{always} and primal feasibility only \emph{at the end} (when $M$ is perfect); it engineers the run so that the moment primal feasibility is reached, CS already holds because $M$ uses only tight edges. This is the design principle of \emph{primal--dual} algorithms: keep one side feasible, drive the other toward feasibility while never violating complementary slackness, and you arrive at a matched optimal pair.
\end{whybox}

\paragraph{Egerv\'ary's theorem as strong duality.}
Egerv\'ary's 1931 theorem is precisely strong duality for this LP, proved combinatorially two decades before linear programming had a name. K\H{o}nig's theorem (max matching $=$ min vertex cover in bipartite graphs) is the special case where all costs are $0/1$: the potentials become a $0/1$ vertex cover. The Hungarian algorithm thus generalizes K\H{o}nig's min-cover / max-matching duality to arbitrary edge weights, with potentials playing the role of a fractional, real-valued cover.

\section{Applications}
\label{ch16:apps}

\begin{applicationbox}
\textbf{Job assignment / personnel scheduling.} The textbook use: $c_{ij}=$ time (or cost, or $-$profit) for worker $i$ on task $j$; minimize the total. Because a perfect matching uses each worker and each task once, the Hungarian algorithm returns the globally cheapest one-to-one assignment in $\Oh(n^3)$. It generalizes to \emph{bottleneck assignment} (minimize the largest edge used) by binary-searching a threshold and testing for a perfect matching among edges below it (Chapter~15).
\end{applicationbox}

\begin{applicationbox}
\textbf{Tracking and data association.} In multi-object tracking, each frame's detections must be matched to existing tracks; $c_{ij}$ is a distance (Mahalanobis, IoU, appearance feature). The Hungarian algorithm produces the maximum-likelihood association each frame, and it is the core of classic trackers (the assignment step inside SORT and Kalman-filter trackers). Gating sets impossible matches to $+\infty$.
\end{applicationbox}

\begin{applicationbox}
\textbf{Optimal transport, discretely.} When two distributions are supported on $n$ points each with equal masses, the optimal-transport (Earth Mover) plan between them is exactly a min-cost \emph{assignment} with $c_{ij}=$ ground distance. The dual potentials $u,v$ are the \emph{Kantorovich potentials}. Sinkhorn's entropic regularization is the fast approximate cousin; the Hungarian algorithm is the exact solver for the unregularized, balanced, equal-mass case.
\end{applicationbox}

\begin{applicationbox}
\textbf{Frequency / channel assignment, antenna pointing, taxi dispatch, wiring placement.} Any setting with two finite sides, a one-to-one constraint, and an additive objective is in scope. When the sides are large but most pairs are forbidden, model as min-cost flow on the sparse graph instead of running dense Hungarian.
\end{applicationbox}

\begin{connectionbox}
The dual potentials reappear across this book. In Chapter~14 they are min-cost-flow node potentials; in Chapter~35 reduced costs play a voltage-like role; and the LP-duality template here -- maintain dual feasibility, enforce complementary slackness -- is the same template behind the primal--dual approximation algorithms of Chapter~38. Total unimodularity, the reason the LP is integral, is the bipartite specialization of the integrality you will meet when matchings stop being bipartite (Chapter~17).
\end{connectionbox}

\begin{aestheticbox}
Even where faster scaling solvers exist for huge sparse graphs, the Hungarian algorithm remains the clearest setting in which the primal--dual method is completely visible: a dual object you can write on a napkin (one number per vertex), a feasibility rule with a one-line proof, a search that the dual itself steers, and a stopping test that doubles as a certificate. Once you understand it, min-cost flow, weighted matching in general graphs, and primal--dual approximation algorithms are variations on the same method.

\textbf{Looking ahead:} Every technique in this chapter relies on the two-sided geometry of bipartite graphs: alternating paths cannot loop back on themselves in odd cycles, and total unimodularity guarantees integral vertices for free. In general graphs, where any vertex can pair with any other, these guarantees fail. In Chapter~17, odd cycles -- which Jack Edmonds called \emph{blossoms} -- break the two-coloring of alternating trees, requiring a cycle-contraction step to restore alternating search.
\end{aestheticbox}

\section{Exercises}
\label{ch16:exercises}

\begin{exercise}[Quick mental check: initial feasibility]
Suppose all edge costs in $C$ are strictly positive ($c_{ij} > 0$). Is the all-zero potential assignment $u_i = 0$ for all $i\in L$ and $v_j = 0$ for all $j\in R$ feasible? Can the Hungarian algorithm begin with these potentials, and if it does, what is the initial equality subgraph?
\end{exercise}

\begin{exercise}[Quick mental check: zero-cost shift]
If you subtract the row minimum $\min_j c_{ij}$ from every row and then the column minimum from every column, explain in one sentence why every resulting entry is non-negative and at least one entry in each row and column is zero.
\end{exercise}

\begin{exercise}
Run the Hungarian algorithm by hand on
\(
\left[\begin{smallmatrix} 9&11&14\\ 6&15&13\\ 12&13&6\end{smallmatrix}\right]
\)
(minimize). Show every row/column reduction, each dual adjustment with its $\delta$ and the sets $S,T$, the augmenting path, and the final $u,v$. Verify that $\sum u_i+\sum v_j$ equals the matching cost.
\end{exercise}

\begin{exercise}
Prove that adding a constant to an entire row (or column) of the cost matrix does not change which assignment is optimal, only its value. Use this to explain why initial row/column reduction is legitimate, and why max-weight and min-cost assignment are interconvertible.
\end{exercise}

\begin{exercise}
\textbf{(Break it.)} Construct a $3\times3$ instance where the greedy rule ``repeatedly pick the cheapest unused cell whose row and column are still free'' returns a strictly suboptimal assignment. Quantify the gap, and identify the alternating cycle the Hungarian algorithm uses to repair greedy's mistake.
\end{exercise}

\begin{exercise}
Implement the $\Oh(n^3)$ algorithm of Section~\ref{ch16:algorithm}. Test it against a brute-force $\Oh(n!)$ checker for all $n\le 8$ on random integer matrices, asserting that (a) the costs agree and (b) your returned $(u,v)$ are feasible and sum to the cost. Then stress-test $n=500$ and report the runtime.
\end{exercise}

\begin{exercise}
Show how to handle a \emph{rectangular} instance with $\abs{L}=p<\abs{R}=q$ (assign every left vertex, each right vertex at most once) by padding to a square. Prove your padding does not change the optimal assignment of the real vertices. What costs do the dummy rows need?
\end{exercise}

\begin{exercise}
\textbf{(Min-cost flow reduction.)} Write out the flow network of the \emph{Connection} box in Section~\ref{ch16:complexity} for the worked example of Section~\ref{ch16:worked}. Run successive-shortest-paths with Johnson potentials and check that the node potentials you obtain match $u,v$ up to the standard sign/offset convention.
\end{exercise}

\begin{exercise}
\textbf{(Total unimodularity.)} Prove that the constraint matrix of the assignment LP is totally unimodular. (Hint: every square submatrix can be expanded along a row or column with at most two $\pm1$ entries; induct, using bipartiteness to two-colour the rows.) Conclude that the LP relaxation has an integral optimum.
\end{exercise}

\begin{exercise}
\textbf{(Sensitivity.)} After solving an instance optimally with potentials $(u,v)$, suppose one cost $c_{i^\star j^\star}$ decreases by $t>0$. Give an $\Oh(n^2)$ procedure to update the optimal matching (hint: at most one augmenting search is needed). For which $t$ does the current matching stay optimal? Express the threshold using the slack of $(i^\star,j^\star)$.
\end{exercise}

\begin{exercise}
\textbf{(Bottleneck assignment.)} You want the perfect matching minimizing the \emph{maximum} edge cost used, not the sum. Give an $\Oh(n^2 \log n)$ algorithm by binary-searching the threshold and testing for a perfect matching among edges below it (using unweighted bipartite matching, Chapter~15). Prove correctness.
\end{exercise}

\begin{exercise}
\textbf{(Stability vs.\ optimality.)} Contrast min-cost assignment with a \emph{stable} matching (Gale--Shapley) on preference-derived costs. Give a small instance whose min-cost matching is not stable, and one whose stable matching is not min-cost. Explain in one paragraph why the two objectives optimize different things.
\end{exercise}

\begin{exercise}
\textbf{(Counting dual steps.)} For the $3\times3$ worked example of Section~\ref{ch16:worked}, exactly one dual adjustment occurred. Construct, for each $k\in\{0,1,2\}$, a $3\times3$ instance whose run uses exactly $k$ dual adjustments in total across all phases. Then argue that on an $n\times n$ instance the total number of dual adjustments is at most $n(n-1)/2$, and exhibit a family approaching that bound.
\end{exercise}

\begin{exercise}
\textbf{(Certificate checking.)} You are handed a claimed solution: an assignment $\pi$ and potentials $(u,v)$. Write the $\Oh(n^2)$ checker that decides whether the pair certifies optimality, and state precisely the three conditions it must verify (feasibility, tightness of matched edges, value agreement). Explain why a verifier needs no access to the algorithm that produced the answer -- this is the practical payoff of carrying a dual.
\end{exercise}
